\documentclass[12pt,a4paper]{article}
\usepackage{ugmath}
\usepackage{float}
\usepackage{placeins}
\usepackage[labelfont=bf,labelsep=space]{caption}
\newcommand{\paren}[1]{\left( #1 \right)}
\newcommand{\alumno}{Ricardo León Martínez}
\newcommand{\materia}{Fisica I}
\newcommand{\profesor}{Lucero Uscanga Aguilera}
\newcommand{\tarea}{Examen 1}
\newcommand{\fecha}{17/3/2026}

\begin{document}
    \begin{center}
    {\large\textbf{UNIVERSIDAD DE GUANAJUATO}}\\[0.3cm]
    {\normalsize\textbf{DIVISIÓN DE CIENCIAS NATURALES Y EXACTAS}}\\
    {\normalsize\textbf{CAMPUS GUANAJUATO}}\\[1cm]

    {\Large\textbf{\tarea\ (\materia)}}\\[1cm]
    \end{center}

    \textbf{Nombre:} \alumno \hfill 
    \textbf{Fecha:} \fecha \hfill 
    \textbf{Calificación:} \rule{3cm}{0.4pt} \\[0.3cm]

    \textbf{1.} La aceleración de un cuerpo que se mueve a lo largo de una linea recta esta
    dada por la ecuación $a=-kv^{2}$, donde $k$ es una constante y suponiendo que $t=0$,
    $v=v_{0}$. Encontrar la velocidad y el dezplasamiento en función del tiempo. Encontrar
    tambien $"v"$ en función de $"x"$.

    \begin{solucion}
        Dado que la acelareción es la derivada de la velocidad entonces
        \begin{equation*}
            \frac{dv}{dt}=-kv^{2}.
        \end{equation*}
        Reescribiendo
        \begin{equation*}
            v^{-2}dv=-kdt.
        \end{equation*}
        Integrando a ambos lados
        \begin{equation*}
            \int v^{-2}dv=-k\int dt.
        \end{equation*}
        Así,
        \begin{equation*}
            -\frac{1}{v}=-kt+C.
        \end{equation*}
        Multiplicando por -1 ambos lados
        \begin{equation*}
            \frac{1}{v}=kt+C^{\prime}
        \end{equation*}
        Si $v$ es una función del tiempo y $v_{0}$ suecede cuando $t=0$ entonces
        \begin{equation*}
            \frac{1}{v_{0}}=C^{\prime}.
        \end{equation*}
        Por tanto
        \begin{equation*}
            \frac{1}{v}=kt+\frac{1}{v_{0}}.
        \end{equation*}
        Finalmente
        \begin{equation*}
            v(t)=\frac{v_{0}}{v_{0}kt+1}
        \end{equation*}
        Ahora, dado que la velocidad es la derivada de la posición se tiene
        \begin{equation*}
            \frac{dx}{dt}=\frac{v_{0}}{v_{0}kt+1}.
        \end{equation*}
        por lo tanto
        \begin{equation*}
            dx=\frac{v_{0}}{v_{0}kt+1}dt.
        \end{equation*}
        Integramos ambos lados
        \begin{equation*}
            \int dx=\int\frac{v_{0}}{v_{0}kt+1}dt
        \end{equation*}
        En consecuencia
        \begin{equation*}
            x=\frac{1}{k}\ln(v_{0}kt+1)+C
        \end{equation*}
        Si tomamos $x(0)=x_{0}$ entonces
        \begin{equation*}
            x_{0}=C
        \end{equation*}
        De esta forma
        \begin{equation*}
            x(t)=\frac{1}{k}\ln(v_{0}kt+1)+x_{0}.
        \end{equation*}
        Así hemos encontrado expresiones de velocidad y posición con respecto al tiempo.
        Ahora usamos
        \begin{equation*}
            a=\frac{dv}{dt}=v\frac{dv}{dx}
        \end{equation*}
        Por lo tanto,
        \begin{equation*}
            v\frac{dv}{dx}=-kv^{2}.
        \end{equation*}
        Si $v\neq0$, dividimos por v ambos lados
        \begin{equation*}
            \frac{dv}{dx}=-kv.
        \end{equation*}
        Reescribiendo
        \begin{equation*}
            \frac{1}{v}dv=-kdx.
        \end{equation*}
        Integramos a ambos lados
        \begin{equation*}
            \int\frac{1}{v}dv=-k\int dx.
        \end{equation*}
        por lo tanto
        \begin{equation*}
            \ln(v)=-kx+C
        \end{equation*}
        Exponenciando
        \begin{equation*}
            v=C^{\prime}e^{-kx}
        \end{equation*}
        Tomando $v(0)=v_{0}$ y $x(0)=0$ se tiene
        \begin{equation*}
            C^{\prime}=v_{0}e^{kx_{0}}.
        \end{equation*}
        Así,
        \begin{equation*}
            v(x)=v_{0}e^{(x_{0}-x)k}
        \end{equation*}
    \end{solucion}
\end{document}